% !TeX root = ../main.tex
% ================================================================
%  sections/04_model_rewrite.tex
%
%  Proposed EJOR-style rewrite of the bilevel formulation.
%  Self-contained: \providecommand guards define every new symbol,
%  so this file compiles alongside styles/macros.tex unchanged.
%
%  It does NOT overwrite 04b_leader.tex / 04c_follower.tex.
%  Symbol migration map: see the table at the end of this file.
% ================================================================

% ── New/renamed symbols (guarded; move into styles/macros.tex) ──
\providecommand{\taxstep}{\Delta}                        % was \gridstep (=\theta): CLASHED with return lag \theta_p
\providecommand{\taxgrid}{\mathcal{T}^{\tau}}
\providecommand{\lexle}{\preceq_{\mathrm{lex}}}
\providecommand{\Mtrip}{M^{\mathrm{trip}}}
\providecommand{\Mveh}{M^{\mathrm{veh}}}
\providecommand{\Mrdij}{\setM^{\mathrm{rd}}_{ij}}
\providecommand{\Moutij}{\setM^{\mathrm{out}}_{ij}}

%=================================================================
\subsection{Notational conventions}
\label{sec:model:conventions}

Three conventions keep the formulation compact. First, an overline or
underline always denotes a \emph{bound} on the quantity carrying it
($\underline{A}_i\le A_i\le\overline{A}_i$), and never any other
operation. Second, whenever a relation is indexed by a pair set, we
write its slices as
\begin{equation}
\setG_p\defeq\{g\in\setG:(p,g)\in\Omega^{\mathrm{p}}\},\quad
\setP_g\defeq\{p\in\setP:(p,g)\in\Omega^{\mathrm{p}}\},\quad
\setL_p,\ \setP_{\ell}\ \text{analogously,}
\label{eq:slices}
\end{equation}
\begin{equation}
\Phi_{ij}\defeq\{q\in\setQ:(i,j)\in\setA_q\},\qquad
\Mrdij\defeq\setM_{ij}\cap\setMrd,\qquad
\Moutij\defeq\setM_{ij}\cap\setMout ,
\label{eq:slices2}
\end{equation}
so that every summation ranges over a named set. Third, aggregate node
flows and the network-activity indicator are
\begin{equation}
X^{\mathrm{in}}_{qit}\defeq\!\!\sum_{j:(j,i)\in\setA_q}\sum_{m\in\setM_{ji}}\!\! x_{qjimt},
\qquad
X^{\mathrm{out}}_{qit}\defeq\!\!\sum_{j:(i,j)\in\setA_q}\sum_{m\in\setM_{ij}}\!\! x_{qijmt},
\qquad
\tilde{y}_i\defeq
\begin{cases}
y_i, & i\in\setNo,\\
1,   & i\in\setR,
\end{cases}
\label{eq:flows}
\end{equation}
for all $q\in\setQ$, $i\in\setN$, $t\in\setT$. Finally, the floor area
used at facility $i$ in period $t$ combines a throughput-handling
component (factories and warehouses) with a stock-storage component
(warehouses only; factories hold no inventory, \Cref{sec:model:c1}):
\begin{equation}
A^{\mathrm{u}}_{it}\defeq
\begin{cases}
\displaystyle\sum_{p\in\setP} a^{\mathrm{hd}}_p\,X^{\mathrm{out}}_{pit},
  & i\in\setF,\\[1.0em]
\displaystyle\sum_{p\in\setP}\Bigl(a^{\mathrm{hd}}_p\,X^{\mathrm{in}}_{pit}
  + a^{\mathrm{st}}_p\,I_{pit}\Bigr),
  & i\in\setW,
\end{cases}
\qquad \forall\, i\in\setF\cup\setW,\ t\in\setT .
\label{eq:used_area}
\end{equation}

%=================================================================
\subsection{The bilevel programme}
\label{sec:model:bilevel}

The regulator (leader) announces a carbon-tax index
$\hat{\tau}\in\taxgrid\defeq\{\underline{\hat{\tau}},\ldots,
\overline{\hat{\tau}}\}\subset\Z$, inducing the rate
$\tau=\taxstep\hat{\tau}$ with grid step $\taxstep$
(USD/kg\,CO$_2$e). The supply chain operator (follower) observes
$\hat{\tau}$ and chooses
\begin{equation}
\bz=\bigl(y,\,A,\,\lambda^{\mathrm{p}},\,\lambda^{\mathrm{r}},\,\rho,\,
\sigma,\,\zeta,\,\kappa,\,n,\,N,\,q^{\mathrm{p}},\,q^{\mathrm{r}},\,
I,\,x\bigr)\in\setX ,
\label{eq:zvec}
\end{equation}
where $\setX$ is the mixed-integer set defined by
\eqref{eq:c1}--\eqref{eq:domains} below. Writing $\lexle$ for the
lexicographic order on $\R^2$ --- $(e_1,\hat\tau_1)\lexle
(e_2,\hat\tau_2)$ iff $e_1<e_2$, or $e_1=e_2$ and
$\hat\tau_1\le\hat\tau_2$ --- the complete programme is
\begin{subequations}\label{eq:BL}
\begin{alignat}{2}
\textup{(BL)}\quad
\lexmin_{\hat{\tau},\,\bz}\quad
 & \bigl(\Emis(\bz),\ \hat{\tau}\bigr) &&
\label{eq:BL_obj}\\[0.2em]
\text{s.t.}\quad
 & \hat{\tau}\in\taxgrid , &&
\label{eq:BL_tax}\\[0.2em]
 & \bz\in\Sol{\hat{\tau}}\defeq
   \argmin_{\bz'\in\setX}\;
   \bigl\{\Cost(\bz')+\taxstep\,\hat{\tau}\,\Emis(\bz')\bigr\}. &&
\label{eq:BL_react}
\end{alignat}
\end{subequations}
Minimising jointly over $(\hat\tau,\bz)$ in~\eqref{eq:BL_obj} encodes
the \emph{optimistic} (strong) bilevel convention
\citep{dempe2002}: when the follower's best response is not
unique, the leader resolves the tie in favour of the lowest-emission
response. The second lexicographic criterion encodes a
\emph{least-intervention principle} --- among tax indices attaining the
minimum emissions, the regulator levies the smallest.

\begin{remark}[Why lexicographic and not weighted-sum]
\label{rem:no_scalarization}
The two leader criteria are incommensurable: $\Emis$ is measured in
kg\,CO$_2$e, $\hat\tau$ is a dimensionless grid index. A weighted sum
$\Emis(\bz)+\varepsilon\hat\tau$ is \emph{exactly} equivalent
to~\eqref{eq:BL_obj} for any
$0<\varepsilon<\Delta E_{\min}/(\overline{\hat\tau}-\underline{\hat\tau})$,
where $\Delta E_{\min}$ is the smallest positive gap between distinct
attainable emission levels; but $\Delta E_{\min}$ is not known a priori
and any safe surrogate drives $\varepsilon$ below solver tolerance. We
therefore keep the pre-emptive form and enforce it algorithmically
(\Cref{sec:method:tiebreak}).
\end{remark}

Since the follower's problem is a \MILP for every fixed $\hat\tau$,
\eqref{eq:BL} is NP-hard \citep{moore1990,jeroslow1985}.

%- - - - - - - - - - - - - - - - - - - - - - - - - - - - - - - - -
\subsubsection{Structural properties}
\label{sec:model:structure}

The tax index enters~\eqref{eq:BL_react} only through an objective
coefficient. This observation has three consequences that the solution
method of \Cref{sec:methodology} exploits.

\begin{proposition}[Parameter-free feasible set]\label{prop:X_indep}
$\setX$ does not depend on $\hat{\tau}$. Consequently
$\Sol{\hat\tau}\neq\emptyset$ for every $\hat\tau\in\taxgrid$ if and
only if $\setX\neq\emptyset$, and bilevel feasibility need be verified
once, not once per tax index.
\end{proposition}

\begin{proposition}[Monotone emission response]\label{prop:monotone}
Let $\hat\tau_1<\hat\tau_2$ in $\taxgrid$ and let
$\bz_1\in\Sol{\hat\tau_1}$, $\bz_2\in\Sol{\hat\tau_2}$ be \emph{any}
optimal responses. Then $\Emis(\bz_2)\le\Emis(\bz_1)$. Hence the
optimal-response emission level
$\Emis^{*}(\hat\tau)\defeq\min\{\Emis(\bz):\bz\in\Sol{\hat\tau}\}$ is a
non-increasing step function of $\hat\tau$, and the follower value
function
$\phi(\hat\tau)\defeq\min_{\bz\in\setX}\{\Cost(\bz)+\taxstep\hat\tau\Emis(\bz)\}$
is concave and piecewise linear in $\hat\tau$, with slope
$\taxstep\,\Emis^{*}(\hat\tau)$.
\end{proposition}

\begin{proof}
Optimality of $\bz_1$ at $\hat\tau_1$ and of $\bz_2$ at $\hat\tau_2$ gives
$\Cost_1+\taxstep\hat\tau_1\Emis_1\le\Cost_2+\taxstep\hat\tau_1\Emis_2$
and
$\Cost_2+\taxstep\hat\tau_2\Emis_2\le\Cost_1+\taxstep\hat\tau_2\Emis_1$.
Adding and cancelling $\Cost_1,\Cost_2$ yields
$\taxstep(\hat\tau_2-\hat\tau_1)(\Emis_2-\Emis_1)\le0$; since
$\taxstep>0$ and $\hat\tau_2>\hat\tau_1$, $\Emis_2\le\Emis_1$.
Concavity of $\phi$ is immediate: it is a pointwise minimum of the
affine functions $\hat\tau\mapsto\Cost(\bz)+\taxstep\hat\tau\Emis(\bz)$
over $\bz\in\setX$, and piecewise linearity follows because $\setX$ has
finitely many extreme optimal responses over the finite grid.
\end{proof}

\begin{corollary}[Threshold form of the leader's problem]
\label{cor:threshold}
The leader's optimal emission level is
$\Emis^{*}(\overline{\hat\tau})$, and its optimal tax index is
$\hat\tau^{*}=\min\{\hat\tau\in\taxgrid:
\Emis^{*}(\hat\tau)=\Emis^{*}(\overline{\hat\tau})\}$. The predicate
$\Emis^{*}(\hat\tau)=\Emis^{*}(\overline{\hat\tau})$ is monotone in
$\hat\tau$ by \Cref{prop:monotone}, so $\hat\tau^{*}$ is located by
bisection in $O(\log|\taxgrid|)$ exact follower solves.
\end{corollary}

\Cref{prop:monotone} also explains, as a property of the model rather
than an empirical regularity, why the emission--tax response observed in
\Cref{sec:experiments} is a saturating staircase: the response is
non-increasing and piecewise constant by construction, and it flattens
once no remaining reconfiguration of the network trades cost for
emissions at the prevailing rate.

%=================================================================
\subsection{Leader: the regulator}
\label{sec:model:leader}

The leader's feasible set is the finite grid~\eqref{eq:BL_tax} together
with the bilevel-optimality condition~\eqref{eq:BL_react}. The index
bounds are \emph{derived} from the substantive policy interval
$[\underline{\tau},\overline{\tau}]$ of \Cref{app:data:tax_range} and
the grid step, as
$\underline{\hat\tau}=\underline{\tau}/\taxstep$ and
$\overline{\hat\tau}=\overline{\tau}/\taxstep$; the same interval is
traversed at any resolution by varying $\taxstep$ alone. The index is a
modelling device confined to~\eqref{eq:BL} and to the \DBP of
\Cref{sec:methodology}: all reported results are expressed in the rate
$\tau$ (USD/t\,CO$_2$e).

%=================================================================
\subsection{Follower: the supply chain operator}
\label{sec:model:follower}

\subsubsection{Objective}
\label{sec:model:follower_obj}

Total cost decomposes into network investment, operations, and
transport, $\Cost(\bz)=\Cost^{\mathrm{inv}}(\bz)+
\Cost^{\mathrm{op}}(\bz)+\Cost^{\mathrm{tr}}(\bz)$, with
\begin{align}
\Cost^{\mathrm{inv}}(\bz)
&=\sum_{s\in\setS}F^{\mathrm{s}}_s\,y_s
 +\sum_{i\in\setF\cup\setW}c^{\mathrm{area}}_i\,A_i
 +\sum_{f\in\setF}\Bigl(\sum_{g\in\setG}C^{\mathrm{p}}_g\lambda^{\mathrm{p}}_{fg}
  +\sum_{\ell\in\setL}C^{\mathrm{r}}_{\ell}\lambda^{\mathrm{r}}_{f\ell}\Bigr)
 +\sum_{i\in\setN}\sum_{m\in\setMrd}F^{\mathrm{v}}_m\,N_{im},
\label{eq:cost_inv}\\[0.4em]
\Cost^{\mathrm{op}}(\bz)
&=\abs{\setT}\sum_{h\in\setH}F^{\mathrm{h}}_h\,y_h
 +\sum_{t\in\setT}\sum_{f\in\setF}\Bigl(
   \sum_{g\in\setG}\sum_{p\in\setP_g}c^{\mathrm{p}}_g\,q^{\mathrm{p}}_{pfgt}
  +\sum_{\ell\in\setL}\sum_{p\in\setP_{\ell}}c^{\mathrm{r}}_{\ell}\,q^{\mathrm{r}}_{pf\ell t}\Bigr)
\nonumber\\
&\quad
 +\sum_{t\in\setT}\Bigl(
   \sum_{i\in\setNv}\sum_{p\in\setP}c^{\mathrm{h}}_p\,I_{pit}
  +\sum_{s\in\setS}\sum_{u\in\setU}c^{\mathrm{u}}_{su}\,X^{\mathrm{out}}_{ust}
  +\sum_{r\in\setR}\sum_{k\in\setK}c^{\mathrm{k}}_k\,X^{\mathrm{out}}_{krt}\Bigr),
\label{eq:cost_op}\\[0.4em]
\Cost^{\mathrm{tr}}(\bz)
&=\sum_{t\in\setT}\sum_{(i,j)\in\setA}\Biggl[
   \sum_{m\in\Mrdij}c^{\mathrm{veh}}_m\,r_m\,\delta_{ij}\,n_{ijmt}
  +\sum_{m\in\Moutij}\sum_{q\in\Phi_{ij}}c^{\mathrm{tr}}_m\,\delta_{ij}\,\omega_q\,x_{qijmt}
\nonumber\\
&\hspace{6.2em}
  +\bigl(\xi^{\mathrm{o}}_{ij}+\xi^{\mathrm{d}}_{ij}\bigr)
   \sum_{m\in\setM_{ij}}\sum_{q\in\Phi_{ij}}c^{\mathrm{hub}}_m\,\omega_q\,x_{qijmt}\Biggr],
\label{eq:cost_tr}
\end{align}
where $\xi^{\mathrm{o}}_{ij}=\mathbf{1}[i\in\setH]$ and
$\xi^{\mathrm{d}}_{ij}=\mathbf{1}[j\in\setH]$ flag hub endpoints. No
separate fixed cost is charged for opening a factory or warehouse:
\eqref{eq:area_bounds} forces $A_i\ge\underline{A}_i>0$ whenever
$y_i=1$, so the area term already prices the opening decision. Road
transport is charged per trip and scaled by the empty-return factor
$r_m$; outsourced transport is charged per kilogram-kilometre.

Total emissions are
\begin{align}
\Emis(\bz)
&=\sum_{i\in\setF\cup\setW}e^{\mathrm{o,f}}_i\,A_i
 +\sum_{t\in\setT}\sum_{f\in\setF}\Bigl(
   \sum_{g\in\setG}\sum_{p\in\setP_g}e^{\mathrm{p}}_g\,\omega_p\,q^{\mathrm{p}}_{pfgt}
  +\sum_{\ell\in\setL}\sum_{p\in\setP_{\ell}}e^{\mathrm{r}}_{\ell}\,\omega_p\,q^{\mathrm{r}}_{pf\ell t}\Bigr)
\nonumber\\
&\quad
 +\sum_{t\in\setT}\sum_{s\in\setS}\sum_{u\in\setU}e^{\mathrm{u}}_{su}\,\omega_u\,X^{\mathrm{out}}_{ust}
 +\sum_{t\in\setT}\sum_{(i,j)\in\setA}\sum_{m\in\setM_{ij}}\sum_{q\in\Phi_{ij}}
   e^{\mathrm{tr}}_m\,r_m\,\omega_q\,\delta_{ij}\,x_{qijmt}.
\label{eq:emissions}
\end{align}
Suppliers and hubs carry no opening emission: a supplier's footprint is
priced entirely through $e^{\mathrm{u}}_{su}$, and hubs are outsourced
terminal-access points (\Cref{as:hub}). The empty return leg of a road
delivery emits although it carries no cargo, hence the same $r_m$ as
in~\eqref{eq:cost_tr}, with $r_m=1$ for outsourced modes.

%- - - - - - - - - - - - - - - - - - - - - - - - - - - - - - - - -
\subsubsection{(C1) Flow conservation and inventory balance}
\label{sec:model:c1}

\begin{subequations}\label{eq:c1}
\begin{alignat}{2}
X^{\mathrm{in}}_{pft}
  +\!\!\sum_{g\in\setG_p}\!\! q^{\mathrm{p}}_{pfgt}
  +\!\!\sum_{\ell\in\setL_p}\!\! q^{\mathrm{r}}_{pf\ell t}
  &= X^{\mathrm{out}}_{pft},
&\quad& \forall\, p\in\setP,\ f\in\setF,\ t\in\setT,
\label{eq:bal_factory}\\
I_{pwt} &= I_{pw,t-1}+X^{\mathrm{in}}_{pwt}-X^{\mathrm{out}}_{pwt},
&& \forall\, p\in\setP,\ w\in\setW,\ t\in\setT,
\label{eq:bal_warehouse}\\
X^{\mathrm{in}}_{kwt} &= X^{\mathrm{out}}_{kwt},
&& \forall\, k\in\setK,\ w\in\setW,\ t\in\setT,
\label{eq:bal_warehouse_k}\\
X^{\mathrm{in}}_{prt} &= d_{rpt},
&& \forall\, p\in\setP,\ r\in\setR,\ t\in\setT,
\label{eq:bal_market}\\
X^{\mathrm{in}}_{qht} &= X^{\mathrm{out}}_{qht},
&& \forall\, q\in\setQ,\ h\in\setH,\ t\in\setT,
\label{eq:bal_hub}\\
I_{pi0} &= I^{0}_{pi},
&& \forall\, p\in\setP,\ i\in\setNv .
\label{eq:inv_init}
\end{alignat}
\end{subequations}
Only warehouses hold stock, and only of finished products; factories,
markets, and hubs are pure flow-through nodes, as are recoverable items
everywhere.

%- - - - - - - - - - - - - - - - - - - - - - - - - - - - - - - - -
\subsubsection{(C2) Reverse logistics and returns}
\label{sec:model:c2}

With $d_{rp,t-\theta_p}\defeq0$ for $t-\theta_p<1$:
\begin{subequations}\label{eq:c2}
\begin{alignat}{2}
\sum_{p\in\setP}\gamma_{kp}\,\underline{\eta}_{rpt}\,d_{rp,t-\theta_p}
  \;\le\; X^{\mathrm{out}}_{krt}
  \;&\le\;\sum_{p\in\setP}\gamma_{kp}\,\overline{\eta}_{rpt}\,d_{rp,t-\theta_p},
&\quad& \forall\, k\in\setK,\ r\in\setR,\ t\in\setT,
\label{eq:return_bounds}\\
X^{\mathrm{in}}_{kft}
  &=\sum_{\ell\in\setL}\sum_{p\in\setP_{\ell}}\beta_{kp}\,q^{\mathrm{r}}_{pf\ell t},
&& \forall\, k\in\setK,\ f\in\setF,\ t\in\setT.
\label{eq:return_factory}
\end{alignat}
\end{subequations}
Every recovered unit entering a factory is consumed by
remanufacturing, closing the material loop.

%- - - - - - - - - - - - - - - - - - - - - - - - - - - - - - - - -
\subsubsection{(C3) Raw material supply}
\label{sec:model:c3}

\begin{subequations}\label{eq:c3}
\begin{alignat}{2}
\underline{S}_{su}\,\sigma_{sut}\;\le\; X^{\mathrm{out}}_{ust}
  \;&\le\;\overline{S}_{su}\,\sigma_{sut},
&\quad& \forall\, s\in\setS,\ u\in\setU,\ t\in\setT,
\label{eq:supply_bounds}\\
\sigma_{sut} &\le y_s,
&& \forall\, s\in\setS,\ u\in\setU,\ t\in\setT,
\label{eq:supplier_link}\\
X^{\mathrm{in}}_{uft}
  &=\sum_{g\in\setG}\sum_{p\in\setP_g}\alpha_{upg}\,q^{\mathrm{p}}_{pfgt},
&& \forall\, u\in\setU,\ f\in\setF,\ t\in\setT.
\label{eq:raw_consumption}
\end{alignat}
\end{subequations}

%- - - - - - - - - - - - - - - - - - - - - - - - - - - - - - - - -
\subsubsection{(C4) Facility opening and area capacity}
\label{sec:model:c4}

\begin{subequations}\label{eq:c4}
\begin{alignat}{2}
\underline{A}_i\,y_i \;\le\; A_i \;&\le\; \overline{A}_i\,y_i,
&\quad& \forall\, i\in\setF\cup\setW,
\label{eq:area_bounds}\\
A^{\mathrm{u}}_{it} &\le A_i,
&& \forall\, i\in\setF\cup\setW,\ t\in\setT,
\label{eq:capacity_use}\\
\underline{I}_{pi}\,y_i \;\le\; I_{pit} \;&\le\;\overline{I}_{pi}\,y_i,
&& \forall\, p\in\setP,\ i\in\setNv,\ t\in\setT.
\label{eq:inv_bounds}
\end{alignat}
\end{subequations}
Installed area is the sole capacity mechanism: it is bounded by
facility-specific limits conditional on opening, must cover the used
area~\eqref{eq:used_area} in every period, and carries the
opening cost and opening emission. Inventory bounds
\eqref{eq:inv_bounds} are folded into this block because they are the
stock-side counterpart of the same open/closed linkage.

%- - - - - - - - - - - - - - - - - - - - - - - - - - - - - - - - -
\subsubsection{(C5) Technology selection and operation}
\label{sec:model:c5}

\begin{subequations}\label{eq:c5}
\begin{alignat}{2}
\underline{Q}^{\mathrm{p}}_g\,\lambda^{\mathrm{p}}_{fg}
  \;\le\;\sum_{p\in\setP_g}q^{\mathrm{p}}_{pfgt}
  \;&\le\;\overline{Q}^{\mathrm{p}}_g\,\lambda^{\mathrm{p}}_{fg},
&\quad& \forall\, f\in\setF,\ g\in\setG,\ t\in\setT,
\label{eq:prod_cap}\\
\underline{Q}^{\mathrm{r}}_{\ell}\,\rho_{f\ell t}
  \;\le\;\sum_{p\in\setP_{\ell}}q^{\mathrm{r}}_{pf\ell t}
  \;&\le\;\overline{Q}^{\mathrm{r}}_{\ell}\,\rho_{f\ell t},
&& \forall\, f\in\setF,\ \ell\in\setL,\ t\in\setT,
\label{eq:rem_cap}\\
\rho_{f\ell t} &\le\lambda^{\mathrm{r}}_{f\ell},
&& \forall\, f\in\setF,\ \ell\in\setL,\ t\in\setT,
\label{eq:rem_operation_link}\\
\lambda^{\mathrm{p}}_{fg} &\le y_f,
&& \forall\, f\in\setF,\ g\in\setG,
\label{eq:tech_link_p}\\
\lambda^{\mathrm{r}}_{f\ell} &\le y_f,
&& \forall\, f\in\setF,\ \ell\in\setL,
\label{eq:tech_link_r}\\
\sum_{g\in\setG_p}\lambda^{\mathrm{p}}_{fg} &\le 1,
&& \forall\, p\in\setP,\ f\in\setF,
\label{eq:tech_excl_p}\\
\sum_{\ell\in\setL_p}\lambda^{\mathrm{r}}_{f\ell} &\le 1,
&& \forall\, p\in\setP,\ f\in\setF.
\label{eq:tech_excl_r}
\end{alignat}
\end{subequations}
Technologies are installed once, at the start of the horizon
(\Cref{as:fixednet}), and shared across compatible products. The two
families are deliberately \emph{not} unified under a single index: a
production technology is either installed or not, whereas an installed
remanufacturing technology is additionally switched on or off per
period through $\rho_{f\ell t}$.

%- - - - - - - - - - - - - - - - - - - - - - - - - - - - - - - - -
\subsubsection{(C6) Transport mode selection and shipments}
\label{sec:model:c6}

\begin{subequations}\label{eq:c6}
\begin{alignat}{2}
\underline{W}_m\,n_{ijmt}
  \;\le\;\sum_{q\in\Phi_{ij}}\omega_q\,x_{qijmt}
  \;&\le\;\overline{W}_m\,n_{ijmt},
&\quad& \forall\,(i,j)\in\setA,\ m\in\Mrdij,\ t\in\setT,
\label{eq:weight_trip}\\
\zeta_{ijmt}\;\le\; n_{ijmt} \;&\le\;\Mtrip_{ijm}\,\zeta_{ijmt},
&& \forall\,(i,j)\in\setA,\ m\in\Mrdij,\ t\in\setT,
\label{eq:trip_link}\\
\sum_{q\in\Phi_{ij}}\omega_q\,x_{qijmt}
  &\le\overline{B}_m\,\zeta_{ijmt},
&& \forall\,(i,j)\in\setA,\ m\in\Moutij,\ t\in\setT,
\label{eq:out_cap}\\
\sum_{m\in\setM_{ij}}\zeta_{ijmt} &\le 1,
&& \forall\,(i,j)\in\setA,\ t\in\setT.
\label{eq:single_mode}
\end{alignat}
\end{subequations}
$\Mtrip_{ijm}$ is an arc- and mode-specific upper bound on the trip
count; \Cref{app:impl_notes:bigm} gives the tightening used in the
implementation.

%- - - - - - - - - - - - - - - - - - - - - - - - - - - - - - - - -
\subsubsection{(C7) Fleet management and node activation}
\label{sec:model:c7}

\begin{subequations}\label{eq:c7}
\begin{alignat}{2}
\zeta_{ijmt} &\le\tilde{y}_i,
&\quad& \forall\,(i,j)\in\setA,\ m\in\setM_{ij},\ t\in\setT,
\label{eq:endpoint_o}\\
\zeta_{ijmt} &\le\tilde{y}_j,
&& \forall\,(i,j)\in\setA,\ m\in\setM_{ij},\ t\in\setT,
\label{eq:endpoint_d}\\
\zeta_{ijmt} &\le\kappa_{mt},
&& \forall\,(i,j)\in\setA,\ m\in\Moutij,\ t\in\setT,
\label{eq:contract_link}\\
\sum_{j:(i,j)\in\setA,\,m\in\setM_{ij}}
  \frac{r_m\,\delta_{ij}}{v_m}\,n_{ijmt}
  &\le H_m\,N_{im},
&& \forall\, i\in\setN,\ m\in\setMrd,\ t\in\setT,
\label{eq:vehicle_hours}\\
N_{im} &\le\Mveh_{im}\,\tilde{y}_i,
&& \forall\, i\in\setN,\ m\in\setMrd,
\label{eq:vehicle_act}\\
N_{im} &\le\Mveh_{im}\!\!\!\sum_{j:(i,j)\in\setA,\,m\in\setM_{ij}}\sum_{t\in\setT}\!\!\! n_{ijmt},
&& \forall\, i\in\setN,\ m\in\setMrd,
\label{eq:vehicle_flow_link}\\
\sum_{i\in\setN}\sum_{m\in\setMrd}F^{\mathrm{v}}_m\,N_{im} &\le\overline{V}, &&
\label{eq:fleet_budget}\\
\sum_{q\in\setQ}\sum_{t\in\setT}
  \bigl(X^{\mathrm{in}}_{qit}+X^{\mathrm{out}}_{qit}\bigr)
  &\ge y_i,
&& \forall\, i\in\setNo.
\label{eq:min_flow}
\end{alignat}
\end{subequations}
Constraint~\eqref{eq:vehicle_hours} presumes $v_m>0$ and $H_m>0$
(\Cref{app:impl_notes:degenerate}); $H_m$ is the driving-hour allowance
per vehicle per period. Constraint~\eqref{eq:min_flow} requires every
activated node to carry incident flow of some item in some period,
ruling out degenerate open-but-unused nodes.

%- - - - - - - - - - - - - - - - - - - - - - - - - - - - - - - - -
\subsubsection{(C8) Intermodal hub routing}
\label{sec:model:c8}

Let $\nu\in\{\mathrm{a},\mathrm{s}\}$ index the hub family (air, sea),
with hub set $\setH^{\nu}$, mode set $\setM^{\nu}$, and per-arc slices
$\setM^{\nu}_{ij}\defeq\setM_{ij}\cap\setM^{\nu}$ and
$\setM^{\neg\nu}_{ij}\defeq\setM_{ij}\setminus\setM^{\nu}$. Hubs are
mode-switching cross-docking points: access volume arriving by
non-$\nu$ modes from non-$\nu$ nodes equals transfer volume departing
to same-family hubs by $\nu$-modes, and symmetrically,
\begin{subequations}\label{eq:c8}
\begin{alignat}{2}
\sum_{i:(i,h)\in\setA_q,\,i\notin\setH^{\nu}}\ \sum_{m\in\setM^{\neg\nu}_{ih}} x_{qihmt}
  &=\sum_{j:(h,j)\in\setA_q,\,j\in\setH^{\nu}}\ \sum_{m\in\setM^{\nu}_{hj}} x_{qhjmt},
&\quad& \forall\,\nu,\ q\in\setQ,\ h\in\setH^{\nu},\ t\in\setT,
\label{eq:hub_ingress}\\
\sum_{i:(i,h)\in\setA_q,\,i\in\setH^{\nu}}\ \sum_{m\in\setM^{\nu}_{ih}} x_{qihmt}
  &=\sum_{j:(h,j)\in\setA_q,\,j\notin\setH^{\nu}}\ \sum_{m\in\setM^{\neg\nu}_{hj}} x_{qhjmt},
&& \forall\,\nu,\ q\in\setQ,\ h\in\setH^{\nu},\ t\in\setT.
\label{eq:hub_egress}
\end{alignat}
\end{subequations}

%- - - - - - - - - - - - - - - - - - - - - - - - - - - - - - - - -
\subsubsection{(C9) Valid inequality}
\label{sec:model:c9}

\begin{equation}
y_f \;\le\;\sum_{g\in\setG}\lambda^{\mathrm{p}}_{fg}
          +\sum_{\ell\in\setL}\lambda^{\mathrm{r}}_{f\ell},
\qquad \forall\, f\in\setF .
\label{eq:vi_tech}
\end{equation}

\begin{remark}[Validity and purpose]\label{rem:valid_ineq}
Under the greenfield convention used throughout
\Cref{sec:experiments}, no arc routes finished product \emph{into} a
factory, so $X^{\mathrm{in}}_{pft}\equiv0$ and \eqref{eq:bal_factory}
collapses to
$\sum_{g\in\setG_p}q^{\mathrm{p}}_{pfgt}+\sum_{\ell\in\setL_p}q^{\mathrm{r}}_{pf\ell t}
=X^{\mathrm{out}}_{pft}$. A factory installing no technology therefore
carries zero outbound flow, and \eqref{eq:raw_consumption} and
\eqref{eq:return_factory} pin its inbound flows to zero as well,
contradicting \eqref{eq:min_flow} if it is open. Hence
\eqref{eq:vi_tech} is implied. Note the disjunction is essential: a
factory engaged purely in remanufacturing satisfies
\eqref{eq:min_flow} through $q^{\mathrm{r}}$ alone, so the narrower
inequality $y_f\le\sum_g\lambda^{\mathrm{p}}_{fg}$ is \emph{not} valid.
\Cref{eq:vi_tech} is retained because it is not redundant in the \LP
relaxation and tightens the bound the \DBP relies on. A second
candidate --- aggregate installed capacity covering aggregate demand
in every period --- is invalid under inventory banking and is
deliberately excluded (\Cref{app:impl_notes:capdemand}).
\end{remark}

%- - - - - - - - - - - - - - - - - - - - - - - - - - - - - - - - -
\subsubsection{Variable domains}
\label{sec:model:domains}

\begin{subequations}\label{eq:domains}
\begin{alignat}{2}
& y_i\in\bin, &\quad& \forall\, i\in\setNo,
\nonumber\\
& \lambda^{\mathrm{p}}_{fg},\ \lambda^{\mathrm{r}}_{f\ell}\in\bin,
  && \forall\, f\in\setF,\ g\in\setG,\ \ell\in\setL,
\nonumber\\
& \rho_{f\ell t}\in\bin, && \forall\, f\in\setF,\ \ell\in\setL,\ t\in\setT,
\nonumber\\
& \sigma_{sut}\in\bin, && \forall\, s\in\setS,\ u\in\setU,\ t\in\setT,
\nonumber\\
& \kappa_{mt}\in\bin, && \forall\, m\in\setMout,\ t\in\setT,
\nonumber\\
& \zeta_{ijmt}\in\bin, && \forall\, (i,j)\in\setA,\ m\in\setM_{ij},\ t\in\setT,
\label{eq:dom_bin}\\[0.35em]
& n_{ijmt}\in\Zplus, && \forall\, (i,j)\in\setA,\ m\in\Mrdij,\ t\in\setT,
\nonumber\\
& N_{im}\in\Zplus, && \forall\, i\in\setN,\ m\in\setMrd,
\label{eq:dom_int}\\[0.35em]
& A_i\in\Rplus, && \forall\, i\in\setF\cup\setW,
\nonumber\\
& q^{\mathrm{p}}_{pfgt}\in\Rplus, && \forall\, (p,g)\in\Omega^{\mathrm{p}},\ f\in\setF,\ t\in\setT,
\nonumber\\
& q^{\mathrm{r}}_{pf\ell t}\in\Rplus, && \forall\, (p,\ell)\in\Omega^{\mathrm{r}},\ f\in\setF,\ t\in\setT,
\nonumber\\
& I_{pit}\in\Rplus, && \forall\, p\in\setP,\ i\in\setNv,\ t\in\setT\cup\{0\},
\nonumber\\
& x_{qijmt}\in\Rplus, && \forall\, (i,j)\in\setA,\ q\in\Phi_{ij},\ m\in\setM_{ij},\ t\in\setT.
\label{eq:dom_cont}
\end{alignat}
\end{subequations}

Constraints \eqref{eq:c1}--\eqref{eq:domains} define $\setX$. Together
with \eqref{eq:cost_inv}--\eqref{eq:emissions} they instantiate the
follower's problem~\eqref{eq:BL_react}, completing the \BLMIP solved in
\Cref{sec:methodology}.

% ================================================================
%  SYMBOL MIGRATION MAP  (old -> new; delete before submission)
% ----------------------------------------------------------------
%   \gridstep = \theta      ->  \taxstep = \Delta   [MANDATORY: clashed with \theta_p]
%   YC_i                    ->  A_i                 [multi-letter italic]
%   YCT_{it}                ->  A^{u}_{it}          [multi-letter italic]
%   apu_q                   ->  a^{st}_p            [multi-letter italic]
%   apur_p                  ->  a^{hd}_p            [multi-letter italic]
%   s^{inv}_{pit}           ->  I_{pit}             [s clashed with supplier index]
%   a_{f l t}               ->  \rho_{f l t}        [frees 'a' for area coefficients]
%   Q_{ijmt}                ->  n_{ijmt}            [Q clashed with \bar Q^p_g and \setQ]
%   K_{im}                  ->  N_{im}              [K clashed with \setK / index k]
%   F_i (supplier)          ->  F^{s}_s
%   F_h (hub)               ->  F^{h}_h
%   F^{m}_m (vehicle)       ->  F^{v}_m
%   \bar y_i                ->  \tilde y_i          [overline reserved for bounds]
%   \bar\setM^{nu}_{ij}     ->  \setM^{nu}_{ij}
%   \setM^{-nu}_{ij}        ->  \setM^{\neg nu}_{ij}
%   M^{trip}_m (vehicles)   ->  M^{veh}_{im}        [MANDATORY: unit mismatch]
%   C1..C10                 ->  C1..C9  [old C5 folded into C4]
% ================================================================
